\documentclass[12pt,a4paper]{article}

\usepackage{fontspec}
\setmainfont{Times New Roman}
\usepackage[vietnamese]{babel}

\usepackage[
top=2.5cm,
bottom=2.5cm,
left=2.5cm,
right=2.5cm
]{geometry}

\usepackage{longtable}
\usepackage{array}
\usepackage{enumitem}
\usepackage{hyperref}

\setlength{\parindent}{0pt}
\setlength{\parskip}{6pt}
\renewcommand{\arraystretch}{1.25}

\begin{document}
	
	\begin{center}
		{\Large\textbf{BẢNG KIỂM SOÁT TỔNG HỢP TÀI LIỆU}}\\[0.3cm]
		{\large\textbf{KAN-12 -- Sprint 1}}
	\end{center}
	
	\section{Quy ước tác nhân}
	
	Toàn bộ tài liệu sử dụng thống nhất bốn tên tác nhân sau:
	
	\begin{itemize}
		\item \textbf{Administrator}: Quản trị viên hệ thống.
		\item \textbf{Teacher}: Giáo viên hoặc người quản lý học thuật.
		\item \textbf{Student}: Học sinh sử dụng hệ thống học tập và luyện thi.
		\item \textbf{Parent}: Phụ huynh theo dõi quá trình học tập của học sinh.
	\end{itemize}
	
	Không sử dụng lẫn lộn các tên như Admin, Lecturer, Learner, Pupil
	hoặc Guardian nếu chưa được nhóm thống nhất thành tác nhân riêng.
	
	\section{Danh sách Use Case thống nhất}
	
	\begin{longtable}{
			|>{\centering\arraybackslash}p{1.6cm}
			|p{8.3cm}
			|p{4.3cm}|
		}
		\hline
		\textbf{Mã} &
		\textbf{Tên Use Case} &
		\textbf{Người phụ trách} \\
		\hline
		\endfirsthead
		
		\hline
		\textbf{Mã} &
		\textbf{Tên Use Case} &
		\textbf{Người phụ trách} \\
		\hline
		\endhead
		
		UC-01 & Đăng nhập hệ thống & Phát Nguyễn Đức \\
		\hline
		
		UC-02 & Quản lý người dùng và phân quyền & Phát Nguyễn Đức \\
		\hline
		
		UC-03 & Quản lý môn học và khung năng lực & Thiện Độ \\
		\hline
		
		UC-04 & Quản lý tài liệu học tập & Thiện Độ \\
		\hline
		
		UC-05 & Quản lý ngân hàng câu hỏi & Thiện Độ \\
		\hline
		
		UC-06 & Quản lý đề kiểm tra và đề thi thử & Thiện Độ \\
		\hline
		
		UC-07 & Thực hiện bài đánh giá năng lực & Thành viên N \\
		\hline
		
		UC-08 & Chấm điểm và phân tích kết quả làm bài & Nguyễn Thanh Thiện \\
		\hline
		
		UC-09 & Tạo hoặc cập nhật hồ sơ năng lực học sinh & Nguyễn Thanh Thiện \\
		\hline
		
		UC-10 & Thiết lập điểm mục tiêu & Thành viên N \\
		\hline
		
		UC-11 & Tạo hoặc cập nhật lộ trình học cá nhân hóa & Nguyễn Thanh Thiện \\
		\hline
		
		UC-12 & Luyện tập thích ứng theo năng lực & Nguyễn Thanh Thiện \\
		\hline
		
		UC-13 & Làm đề thi thử & Thành viên N \\
		\hline
		
		UC-14 & Tương tác với AI Tutor & Nguyễn Thanh Thiện \\
		\hline
		
		UC-15 & Theo dõi tiến độ học tập & Thiện Độ \\
		\hline
		
		UC-16 & Xem báo cáo năng lực & Thiện Độ \\
		\hline
		
		UC-17 & Dự đoán khả năng đạt điểm mục tiêu & Nguyễn Thanh Thiện \\
		\hline
		
		UC-18 & Nhận thông báo và cảnh báo & Thành viên N \\
		\hline
		
		UC-19 & Cấu hình và giám sát hệ thống & Phát Nguyễn Đức \\
		\hline
		
	\end{longtable}
	
	\section{Kiểm tra phần Nguyễn Thanh Thiện}
	
	\begin{longtable}{
			|p{8.8cm}
			|>{\centering\arraybackslash}p{2.3cm}
			|p{3.3cm}|
		}
		\hline
		\textbf{Nội dung kiểm tra} &
		\textbf{Trạng thái} &
		\textbf{Ghi chú} \\
		\hline
		\endfirsthead
		
		\hline
		\textbf{Nội dung kiểm tra} &
		\textbf{Trạng thái} &
		\textbf{Ghi chú} \\
		\hline
		\endhead
		
		Mục tiêu tổng quát và mục tiêu cụ thể &
		Hoàn thành &
		File kế hoạch quản lý dự án \\
		\hline
		
		Danh sách rủi ro và phương án xử lý &
		Hoàn thành &
		Đã có bảng rủi ro \\
		\hline
		
		Quy trình Agile Scrum &
		Hoàn thành &
		Đã mô tả quy trình \\
		\hline
		
		Sprint Planning, Daily Scrum, Sprint Review và Retrospective &
		Hoàn thành &
		Đã mô tả từng hoạt động \\
		\hline
		
		Kế hoạch quản lý chất lượng và review &
		Hoàn thành &
		Đã có điều kiện hoàn thành Task \\
		\hline
		
		Kế hoạch đào tạo thành viên &
		Hoàn thành &
		Đã có bảng kế hoạch \\
		\hline
		
		Các giai đoạn và mốc công việc &
		Hoàn thành &
		Đã có kế hoạch Sprint \\
		\hline
		
		Bảng phân công trách nhiệm &
		Hoàn thành &
		Đã sử dụng ma trận RACI \\
		\hline
		
		Kế hoạch giao tiếp và báo cáo tiến độ &
		Hoàn thành &
		Đã có bảng giao tiếp \\
		\hline
		
		Quản lý tài liệu LaTeX &
		Hoàn thành &
		Đã có quy định cấu trúc file \\
		\hline
		
		Quản lý Jira và GitHub &
		Hoàn thành &
		Đã có quy định sử dụng \\
		\hline
		
		Quy định branch, commit và Pull Request &
		Hoàn thành &
		Đã có ví dụ đặt tên \\
		\hline
		
		Công nghệ, công cụ và hạ tầng &
		Hoàn thành &
		Đã có bảng công nghệ \\
		\hline
		
		UC-08 &
		Hoàn thành &
		Đã phân tích đầy đủ \\
		\hline
		
		UC-09 &
		Hoàn thành &
		Đã phân tích đầy đủ \\
		\hline
		
		UC-11 &
		Hoàn thành &
		Đã phân tích đầy đủ \\
		\hline
		
		UC-12 &
		Hoàn thành &
		Đã phân tích đầy đủ \\
		\hline
		
		UC-14 &
		Hoàn thành &
		Đã phân tích đầy đủ \\
		\hline
		
		UC-17 &
		Hoàn thành &
		Đã phân tích đầy đủ \\
		\hline
		
	\end{longtable}
	
	\section{Theo dõi nội dung của các thành viên}
	
	\begin{longtable}{
			|p{3.2cm}
			|p{6.0cm}
			|p{2.6cm}
			|p{3.0cm}|
		}
		\hline
		\textbf{Thành viên} &
		\textbf{Nội dung cần nhận} &
		\textbf{Trạng thái} &
		\textbf{Pull Request} \\
		\hline
		\endfirsthead
		
		\hline
		\textbf{Thành viên} &
		\textbf{Nội dung cần nhận} &
		\textbf{Trạng thái} &
		\textbf{Pull Request} \\
		\hline
		\endhead
		
		Phát Nguyễn Đức &
		Phạm vi, WBS, tổng quan SRS, tác nhân, Use Case Diagram,
		UC-01, UC-02, UC-19 và ma trận phân quyền &
		Chưa nhận &
		Chưa có \\
		\hline
		
		Thiện Độ &
		Phân tích sản phẩm, UC-03 đến UC-06, UC-15, UC-16,
		Screen Flow Administrator và Teacher &
		Chưa nhận &
		Chưa có \\
		\hline
		
		Thành viên N &
		UC-07, UC-10, UC-13, UC-18,
		Screen Flow Student và Parent &
		Chưa nhận &
		Chưa có \\
		\hline
		
	\end{longtable}
	
	\section{Bảng kiểm tra khi nhận Pull Request}
	
	Đối với mỗi Pull Request của thành viên, kiểm tra lần lượt:
	
	\begin{enumerate}
		\item Branch có được tạo từ \texttt{develop} hay không.
		\item Tên branch và commit có chứa mã Jira hay không.
		\item File có nằm đúng thư mục \texttt{latex/chapters} hay không.
		\item Thành viên có chỉnh sửa nhầm file của người khác hay không.
		\item Có commit file tạm như \texttt{.aux}, \texttt{.log},
		\texttt{.toc}, \texttt{.lof} hoặc \texttt{.lot} hay không.
		\item Nội dung có đúng với Description trên Jira hay không.
		\item Tác nhân có sử dụng đúng quy ước chung hay không.
		\item Mã Use Case có trùng hoặc thiếu hay không.
		\item Tên Use Case có đúng với danh sách thống nhất hay không.
		\item Screen Flow có thể hiện đầy đủ màn hình bắt đầu,
		xử lý và kết quả hay không.
		\item Nội dung LaTeX có biên dịch được hay không.
		\item Bảng, hình và chữ có vượt lề hay không.
	\end{enumerate}
	
	\section{Bảng kiểm tra tích hợp cuối}
	
	\begin{longtable}{
			|p{9.5cm}
			|>{\centering\arraybackslash}p{2.3cm}
			|p{2.5cm}|
		}
		\hline
		\textbf{Nội dung kiểm tra} &
		\textbf{Trạng thái} &
		\textbf{Ghi chú} \\
		\hline
		\endfirsthead
		
		\hline
		\textbf{Nội dung kiểm tra} &
		\textbf{Trạng thái} &
		\textbf{Ghi chú} \\
		\hline
		\endhead
		
		Đã ghép nội dung của tất cả thành viên &
		Chưa thực hiện &
		Chờ file thành viên \\
		\hline
		
		Có đủ UC-01 đến UC-19 &
		Chưa thực hiện &
		Chờ file thành viên \\
		\hline
		
		Không có mã Use Case bị trùng &
		Chưa thực hiện &
		Chờ file thành viên \\
		\hline
		
		Tên Use Case thống nhất trong toàn tài liệu &
		Chưa thực hiện &
		Chờ file thành viên \\
		\hline
		
		Tên tác nhân thống nhất &
		Chưa thực hiện &
		Chờ file thành viên \\
		\hline
		
		Use Case Diagram khớp danh sách Use Case &
		Chưa thực hiện &
		Chờ sơ đồ \\
		\hline
		
		Use Case khớp với Screen Flow &
		Chưa thực hiện &
		Chờ Screen Flow \\
		\hline
		
		Screen Flow khớp với ma trận phân quyền &
		Chưa thực hiện &
		Chờ ma trận quyền \\
		\hline
		
		Mục lục được cập nhật đầy đủ &
		Chưa thực hiện &
		Thực hiện sau tích hợp \\
		\hline
		
		Không còn tham chiếu ``??'' &
		Chưa thực hiện &
		Thực hiện sau tích hợp \\
		\hline
		
		Không có bảng hoặc hình vượt lề &
		Chưa thực hiện &
		Thực hiện sau tích hợp \\
		\hline
		
		Tài liệu tổng hợp biên dịch thành công &
		Chưa thực hiện &
		Thực hiện sau tích hợp \\
		\hline
		
	\end{longtable}
	
\end{document}